% ============================================================
% CAPÍTULO 5 — CONTRIBUIÇÕES DA ENCÍCLICA MAGNIFICA HUMANITAS
%              PARA A REGULAÇÃO DA INTELIGÊNCIA ARTIFICIAL
% ============================================================

\chapter{Contribuições da Magnifica Humanitas para a Regulação da Inteligência Artificial}
\label{cap:contribuicoes}

\section{Introdução ao Capítulo}
\label{sec:intro-cap5}

Os capítulos precedentes estabeleceram as bases deste trabalho: a fundamentação antropológica e ética da \MH{} (capítulo 2), os desafios específicos que a IA impõe à personalidade humana (capítulo 3) e os instrumentos jurídicos brasileiros disponíveis para sua proteção (capítulo 4). Este quinto capítulo articula esses três eixos, propondo uma síntese das contribuições que a \MH{} pode oferecer para o aperfeiçoamento da regulação da inteligência artificial.

A tese central deste capítulo é que a \MH{}, embora não seja um documento jurídico nem proponha regras específicas, oferece uma fundamentação antropológica robusta que pode preencher lacunas significativas do debate regulatório contemporâneo. Essas contribuições organizam-se em três planos: a crítica ao paradigma tecnocrático (\ref{sec:critica-tecnocratico}), a formulação de princípios éticos para a governança algorítmica (\ref{sec:principios-eticos-mh}), e a articulação concreta com os instrumentos jurídicos brasileiros (\ref{sec:articulacao-juridica}).

\section{A Crítica ao Paradigma Tecnocrático}
\label{sec:critica-tecnocratico}

\subsection{O Paradigma Tecnocrático como Risco Sistêmico}
\label{subsec:paradigma-risco-sistemico}

Uma das contribuições mais originais da \MH{} para o debate sobre a regulação da IA é a formulação do conceito de ``paradigma tecnocrático''. Se o capítulo segundo examinou esse conceito em sua dimensão antropológica, importa agora explorar suas implicações para o direito e para a regulação.

O paradigma tecnocrático, como desenvolvido pela \MH{} a partir de uma reflexão já presente na encíclica \textit{Laudato Si'} (§106–114), designa não apenas um conjunto de tecnologias, mas uma visão de mundo que reduz a realidade a seus aspectos mensuráveis e manipuláveis, subordinando todas as dimensões da existência humana à lógica da eficiência técnica e do lucro \cite{Francisco2015LS, LeaoXIV2026MH}.

Para a regulação da IA, a crítica ao paradigma tecnocrático tem consequências diretas. Em primeiro lugar, ela revela que os riscos da IA não são meramente acidentais ou decorrentes de má aplicação, mas podem ser sistêmicos, isto é, inerentes à própria lógica de desenvolvimento tecnológico quando dissociada de uma reflexão ética adequada. Em segundo lugar, ela adverte contra a tentação de resolver problemas éticos com soluções puramente técnicas (a chamada ``solução tecnocrática''), que frequentemente agravam os problemas que pretendem resolver.

Essa perspectiva oferece uma base crítica para avaliar propostas regulatórias que se limitam a estabelecer regras procedimentais sem questionar os pressupostos antropológicos do desenvolvimento tecnológico. Para a \MH{}, a regulação da IA não pode ser reduzida a uma questão de compliance técnica; ela deve envolver uma reflexão substantiva sobre que tipo de sociedade queremos construir e que lugar reservamos à pessoa humana.

\subsection{A Recusa da Neutralidade Tecnológica}
\label{subsec:recusa-neutralidade}

A \MH{} oferece também uma crítica à tese da neutralidade tecnológica — a ideia de que a tecnologia é intrinsecamente neutra e que seus efeitos dependem exclusivamente do uso que se faz dela. Para a Encíclica, essa tese ignora que as tecnologias incorporam valores, pressupostos e visões de mundo que moldam a experiência humana de forma não trivial \cite{LeaoXIV2026MH}.

No campo da regulação da IA, a recusa da neutralidade tecnológica tem implicações importantes. Se os sistemas de IA não são neutros — se eles carregam em seu próprio desenho técnico escolhas valorativas sobre o que é relevante, justo ou desejável —, então a regulação não pode limitar-se a controlar usos abusivos, mas deve também incidir sobre o próprio desenho dos sistemas.

\section{Princípios Éticos para a Governança Algorítmica}
\label{sec:principios-eticos-mh}

\subsection{Dignidade Humana como Critério Regulatório}
\label{subsec:dignidade-criterio}

O primeiro e mais fundamental princípio que a \MH{} oferece para a regulação da IA é a dignidade da pessoa humana como critério substantivo de avaliação e limite intransponível para o desenvolvimento tecnológico.

Aplicada à regulação algorítmica, a dignidade opera em três dimensões complementares: como limite negativo (proibindo usos da IA que reduzam a pessoa a objeto), como dever de proteção (exigindo que o Estado crie mecanismos eficazes de tutela) e como mandato de promoção (orientando políticas de desenvolvimento tecnológico que coloquem a pessoa no centro).

Na prática regulatória, essas três dimensões traduzem-se em exigências concretas: a proibição de sistemas de IA que instrumentalizem a pessoa sem seu consentimento informado; a obrigatoriedade de avaliações de impacto algorítmico prévias ao lançamento de sistemas de alto risco; e a previsão de mecanismos de supervisão humana efetiva que garantam o controle último da pessoa sobre as decisões que a afetam.

\subsection{Transparência e Direito à Explicação}
\label{subsec:transparencia-explicacao}

A \MH{} afirma, em consonância com o §125 que trata do direito à explicação, que ``a pessoa humana tem o direito de compreender os critérios que orientam as decisões que lhe dizem respeito'' \cite{LeaoXIV2026MH}. Esse princípio tem implicações diretas para a regulação da IA, especialmente no que tange às exigências de explicabilidade e transparência.

O direito à explicação, como desenvolvido pela \MH{}, não se confunde com a mera disponibilização de informações técnicas sobre o funcionamento do algoritmo. Ele exige que a explicação seja acessível, compreensível e significativa para a pessoa afetada pela decisão. Isso implica um dever positivo de traduzir o funcionamento técnico do sistema em linguagem acessível e de fornecer os fundamentos concretos da decisão particular.

Esse entendimento alinha-se com as propostas mais avançadas da doutrina especializada, que distinguem entre explicabilidade global (compreensão do funcionamento geral do sistema) e explicabilidade local (compreensão das razões de uma decisão particular) \cite{DoshiVelez2017TMI, Selbst2018INC}. A \MH{} oferece uma fundamentação antropológica para essa dupla exigência: a pessoa precisa compreender tanto a lógica geral dos sistemas que a afetam quanto as razões específicas de cada decisão que recai sobre si.

\subsection{Subsidiariedade e Governança Participativa}
\label{subsec:subsidiariedade-governanca}

O princípio da subsidiariedade, central na doutrina social da Igreja, oferece um critério relevante para a estruturação da governança algorítmica. A \MH{} (§105) recorda que ``nenhuma instância superior deve assumir funções que podem ser adequadamente desempenhadas pelas instâncias inferiores'' \cite{LeaoXIV2026MH}.

Aplicado à regulação da IA, o princípio da subsidiariedade implica que a regulação não deve excluir a responsabilidade dos próprios desenvolvedores e usuários de sistemas de IA. O Estado não deve substituir-se à iniciativa privada nem à sociedade civil na definição de padrões éticos e técnicos, mas deve criar as condições para que esses atores exerçam sua responsabilidade de forma efetiva.

Essa perspectiva sugere um modelo de governança algorítmica que articule diferentes níveis: autorregulação regulada (com padrões éticos definidos pelos próprios desenvolvedores, sob supervisão pública), codesign normativo (com participação de múltiplos stakeholders na definição de regras), e fiscalização pública independente (com sanções proporcionais para infrações).

\subsection{Solidariedade e Inclusão Digital}
\label{subsec:solidariedade-inclusao}

O princípio da solidariedade, também central na doutrina social da Igreja, orienta a regulação da IA para a superação das assimetrias e desigualdades que a tecnologia pode aprofundar. A \MH{} enfatiza que o desenvolvimento tecnológico deve beneficiar toda a família humana e não apenas alguns privilegiados (§189), e que a exclusão digital é uma forma contemporânea de exclusão social que viola a dignidade da pessoa \cite{LeaoXIV2026MH}.

Para a regulação da IA, a solidariedade traduz-se em exigências de inclusão digital, letramento algorítmico e acesso equitativo às tecnologias. Essas dimensões são frequentemente negligenciadas nos debates regulatórios, que tendem a concentrar-se nos riscos sem considerar as oportunidades e a distribuição desigual de benefícios.

\section{Articulação com os Instrumentos Jurídicos Brasileiros}
\label{sec:articulacao-juridica}

\subsection{Adequação e Aperfeiçoamento da LGPD}
\label{subsec:lgpd-aperfeicoamento}

A LGPD, como examinado no capítulo anterior, já contém instrumentos relevantes para a proteção da personalidade contra riscos algorítmicos, especialmente o direito à revisão de decisões automatizadas (art. 20) e a avaliação de impacto à proteção de dados (art. 38). Contudo, a perspectiva oferecida pela \MH{} permite identificar lacunas e propor aperfeiçoamentos.

Em primeiro lugar, o direito à revisão de decisões automatizadas poderia ser fortalecido com a explicitação de que a revisão deve ser realizada por pessoa natural com competência técnica e autonomia para modificar a decisão — e não mera formalidade de chancela automática. A \MH{} fundamenta essa exigência na necessidade de que a decisão final sobre a pessoa seja efetivamente humana.

Em segundo lugar, a avaliação de impacto poderia ser ampliada para incluir não apenas riscos à privacidade e proteção de dados, mas também riscos à dignidade, à não discriminação, à autonomia e ao desenvolvimento da personalidade. Trata-se de transformar a AIPD em uma verdadeira Avaliação de Impacto Algorítmico (AIA), inspirada no princípio de que todos os aspectos da personalidade merecem proteção.

Em terceiro lugar, a LGPD poderia prever mecanismos mais efetivos de transparência algorítmica, incluindo a obrigatoriedade de auditorias periódicas independentes para sistemas de IA de alto risco e a publicação de relatórios de impacto algorítmico em formato acessível.

\subsection{Contribuições para o PL 2.338/2023}
\label{subsec:contribuicoes-pl}

O PL 2.338/2023, atualmente em tramitação, representa uma oportunidade ímpar para incorporar, no direito positivo brasileiro, as contribuições da reflexão ética e antropológica da \MH{}.

Em primeiro lugar, o PL poderia explicitar, em seus fundamentos, a dignidade da pessoa humana como critério nuclear de avaliação dos sistemas de IA. Embora a dignidade já seja fundamento constitucional, sua explicitação no marco regulatório da IA teria o efeito de orientar a interpretação de todas as demais disposições e de sinalizar o compromisso do ordenamento com a proteção substantiva da pessoa — e não apenas com requisitos formais de compliance.

Em segundo lugar, o PL poderia fortalecer as exigências de transparência e supervisão humana para sistemas de IA que afetam direitos fundamentais, inspirando-se na distinção entre os três níveis de transparência propostos pela \MH{}: acesso à informação, compreensão do funcionamento e capacidade de contestação.

Em terceiro lugar, o PL poderia prever mecanismos de governança participativa que envolvam não apenas especialistas e representantes do setor, mas também a sociedade civil e as comunidades potencialmente afetadas pelos sistemas de IA. A subsidiariedade, como princípio de justiça social, exige que aqueles que serão afetados pelas decisões algorítmicas tenham voz nos processos decisórios que as definem.

\section{Síntese do Capítulo}
\label{sec:sumario-cap5}

A \MH{} oferece contribuições significativas para a regulação da IA em três planos. No plano da \textbf{crítica}, fornece instrumentos conceituais para questionar o paradigma tecnocrático e a tese da neutralidade tecnológica. No plano dos \textbf{princípios}, oferece fundamentos antropológicos robustos para a dignidade, a transparência, a subsidiariedade e a solidariedade como critérios regulatórios. No plano da \textbf{articulação concreta}, aponta direções para o aperfeiçoamento da LGPD e para o enriquecimento do PL 2.338/2023.

Essas contribuições não substituem o trabalho técnico da regulação, mas conferem-lhe uma fundamentação ética e antropológica que, como se argumentou ao longo do trabalho, é indispensável para que a regulação da IA não seja reduzida a um exercício de mera gestão de riscos, mas se constitua como verdadeira proteção da pessoa humana no contexto da transformação tecnológica.
